\documentclass[11pt,a4paper]{article}

\usepackage{ctex}
\usepackage[utf8]{inputenc}
\usepackage{graphicx}
\usepackage{subcaption}
\usepackage{multirow}
\usepackage{enumitem}
\usepackage{geometry}
\usepackage{hyperref}
\hypersetup{
  colorlinks=true,  % 启用彩色链接
  linkcolor=blue,   % 内部链接颜色（目录、交叉引用等）
  urlcolor=blue,    % URL链接颜色
  citecolor=blue,   % 引用链接颜色
  bookmarksdepth=3, % PDF书签深度
  pdfstartview=FitH % PDF打开时的视图方式
}

\geometry{margin=2.5cm}

\usepackage{listings}
\usepackage[dvipsnames]{xcolor}
\usepackage[framemethod=TikZ]{mdframed}
\usepackage{tikz}  % 添加TikZ包用于绘图
\usetikzlibrary{shapes,positioning}
% 在导言区进行样式设置
\lstset{
  language=Python, % 设置语言
  basicstyle=\ttfamily\small, % 设置字体族和大小
  breaklines=true, % 自动换行
  keywordstyle=\bfseries\color{Purple}, % C语言关键字：紫色粗体
  morekeywords={uint32_t,uint8_t,float,char,static,const}, % C语言扩展关键字
  emph=[1]{delay_init,key_init,lcd_init,infrared_init,rgb_init,buzzer_init,led_init,usart_init,htu21d_init,fan_pwm_init}, % 初始化函数
  emphstyle=[1]\bfseries\color{Red}, % 初始化函数：红色粗体
  emph=[2]{htu21d_t,htu21d_h,get_infrared_status,fan_pwm_control,rgb_ctrl,led_control,buzzer_tweet}, % 硬件控制函数
  emphstyle=[2]\bfseries\color{Orange}, % 硬件控制函数：橙色粗体
  emph=[3]{update_display,get_centered_x,check_uart_command,send_system_status,auto_control,check_human_presence}, % 逻辑函数
  emphstyle=[3]\bfseries\color{Blue}, % 逻辑函数：蓝色粗体
  emph=[4]{sprintf,printf,strncmp,strlen,delay_ms,delay_count}, % 标准库函数
  emphstyle=[4]\color{Magenta}, % 标准库函数：洋红色
  commentstyle=\itshape\color{ForestGreen}, % 注释：森林绿斜体
  stringstyle=\color{Maroon}, % 字符串：栗色
  identifierstyle=\color{Black}, % 标识符：黑色
  numberstyle=\tiny\color{Gray}, % 数字：灰色小字
  columns=flexible,
  numbers=none, % 不显示行号
  % numbersep=1.5em, % 设置行号的具体位置
  % numberstyle=\footnotesize\color{Gray}, % 行号样式
  frame=single, % 边框
  framesep=1em, % 设置代码与边框的距离
  rulecolor=\color{lightgray}, % 边框颜色
  backgroundcolor=\color{Gray!5}, % 背景色：淡灰色
  tabsize=4, % Tab键宽度
  showstringspaces=false % 不显示字符串中的空格标记
}

\begin{document}

\begin{minipage}[t]{0.7\textwidth}
  \fontsize{20pt}{24pt}\selectfont Python 期末复习模拟卷笔记
\end{minipage}
\begin{minipage}[p]{0.25\textwidth}
  \hfill
  \includegraphics[width=3cm]{./sticker.pdf}
\end{minipage}

\tableofcontents

\newpage

\vspace{1cm}

\raggedright

\section{单选题}

\subsection{第1题}
下面代码输出的结果是（）
\begin{lstlisting}
def func(num):
    num += 1
a = 10
func(a)
print(a)
\end{lstlisting}

\begin{enumerate}[label=\Alph*.]
  \item 10
  \item 11
  \item int
  \item 程序执行错误
\end{enumerate}

\textbf{\color{red}【笔记】}

正确答案：A

Python中函数是传值调用，不影响原本的值。

\subsection{第2题}
切片操作\texttt{list(range(6))[::-1]}执行结果为（）

\begin{enumerate}[label=\Alph*.]
  \item \verb|[|0, 2, 4, 6\verb|]|
  \item \verb|[|6, 5, 4, 3, 2, 1\verb|]|
  \item \verb|[|0, -1\verb|]|
  \item \verb|[|5, 4, 3, 2, 1, 0\verb|]|
\end{enumerate}

\begin{mdframed}[linewidth=1pt, linecolor=black]

  \textbf{\color{red}【笔记】}

  正确答案：D

  \texttt{list()}将对象转换成列表类型，\texttt{[::-1]}将列表反转。

\end{mdframed}

\subsection{第3题}
下列选项中可以获取Python整数类型帮助的是（）

\begin{enumerate}[label=\Alph*.]
  \item help(float)
  \item dir(float)
  \item help(int)
  \item dir(str)
\end{enumerate}

\subsection{第4题}
以下选项中符合Python语言变量命名规则的是（）

\begin{enumerate}[label=\Alph*.]
  \item it's
  \item 3C
  \item pass
  \item \_AI
\end{enumerate}

\begin{mdframed}[linewidth=1pt, linecolor=black]

  \textbf{\color{red}【笔记】}

  正确答案：D

  \begin{itemize}
    \item A选项中出现了\texttt{'}符号，在Python中的意思是字符串；
    \item B选项数字开头是不允许的；
    \item C选项是Python关键字；
    \item D选项符合Python变量名要求。
  \end{itemize}

\end{mdframed}

\subsection{第5题}
关于Python序列类型的通用操作符和函数，以下选项中描述错误的是（）

\begin{enumerate}[label=\Alph*.]
  \item 如果x是s的元素，x in s 返回True
  \item 如果s是一个序列，s = [1,"kate",True]，s[3] 返回True
  \item 如果s是一个序列，s = [1,"kate",True]，s[-1] 返回True
  \item 如果x不是s的元素，x not in s 返回True
\end{enumerate}

\begin{mdframed}[linewidth=1pt, linecolor=black]

  \textbf{\color{red}【笔记】}

  正确答案：B

  列表s的长度为3，可以访问的索引为0,1,2，如果访问s[3]会报错。

\end{mdframed}

\subsection{第6题}
下列选项中不是Python保留字的是（）

\begin{enumerate}[label=\Alph*.]
  \item False
  \item True
  \item do
  \item class
\end{enumerate}

\begin{mdframed}[linewidth=1pt, linecolor=black]

  \textbf{\color{red}【笔记】}

  正确答案：C

  do是C语言关键字，Python中无此关键字。

\end{mdframed}

\subsection{第7题}
Python 3.x语言中，以下表达式输出结果为66的选项是（）

\begin{enumerate}[label=\Alph*.]
  \item print("6+6")
  \item print(6+6)
  \item print(eval("6+6"))
  \item print(eval("6" + "6"))
\end{enumerate}

\begin{mdframed}[linewidth=1pt, linecolor=black]

  \textbf{\color{red}【笔记】}

  正确答案：D

  \begin{itemize}
    \item A选项打印的是字符串"6+6"；
    \item B选项打印的是6+6的值12；
    \item C选项使用eval将两个字符串'6'转换成数字6后相加，得到12并输出；
    \item D选项使用eval将字符串'66'转换成数字66并输出。
  \end{itemize}

\end{mdframed}

\subsection{第8题}
关于eval函数，以下选项中描述错误的是（）

\begin{enumerate}[label=\Alph*.]
  \item eval函数的作用是将输入的字符串转为Python语句，并执行该语句
  \item 如果用户希望输入一个数字，并用程序对这个数字进行计算，可以采用eval(input())组合
  \item 执行eval("Hello")和执行eval("'Hello'")得到相同的结果
  \item 执行eval('123')输出123
\end{enumerate}

\begin{mdframed}[linewidth=1pt, linecolor=black]

  \textbf{\color{red}【笔记】}

  正确答案：C

  \textbf{eval函数详解：}

  \textbf{1. 基本功能：}
  \begin{itemize}
    \item eval()函数用于执行字符串表达式，并返回表达式的结果
    \item 语法：eval(expression, globals=None, locals=None)
    \item expression：要执行的字符串表达式
  \end{itemize}

  \textbf{2. 使用示例：}
    \begin{lstlisting}
# 数学表达式
eval("2 + 3")          # 返回 5
eval("2 * 3 + 1")      # 返回 7
eval("abs(-5)")        # 返回 5

# 字符串表达式
eval("'Hello'")        # 返回 'Hello'
eval("'3' + '4'")      # 返回 '34'

# 变量表达式
x = 10
eval("x + 5")          # 返回 15

# 列表和字典
eval("[1, 2, 3]")      # 返回 [1, 2, 3]
eval("{'a': 1, 'b': 2}")  # 返回 {'a': 1, 'b': 2}
    \end{lstlisting}

  \textbf{3. 错误示例分析：}
  \begin{itemize}
    \item eval("Hello")：会报错，因为Hello被当作变量名，但未定义
    \item eval("'Hello'")：正确，返回字符串'Hello'
    \item eval("123")：返回整数123
    \item eval("'123'")：返回字符串'123'
  \end{itemize}

  \textbf{4. 常见用法：}
    \begin{lstlisting}
# 用户输入处理
user_input = input("请输入表达式: ")
result = eval(user_input)  # 注意：有安全风险

# 配置文件解析
config_str = "{'debug': True, 'port': 8080}"
config = eval(config_str)
    \end{lstlisting}
\end{mdframed}

\subsection{第9题}
下面代码的输出结果是（）
\begin{lstlisting}
d = {"大海":"蓝色", "天空":"灰色", "大地":"黑色"}
print(d["大地"], d.get("大地","黄色"), d.setdefault('草地','绿色'))
\end{lstlisting}

\begin{enumerate}[label=\Alph*.]
  \item 黑色 黑色 None
  \item 黑色 黄色 绿色
  \item 黑色 黄色 None
  \item 黑色 黑色 绿色
\end{enumerate}

\begin{mdframed}[linewidth=1pt, linecolor=black]

  \textbf{\color{red}【笔记】}

  正确答案：D

  \begin{itemize}
    \item 第一个d["大地"]在字典d中存在键，对应为"蓝色"；
    \item 第二个d.get("大地","黄色")，其含义为获取键“大地”，如果不存在返回值“黄色”
    \item 第三个d.setdefault('草地','绿色')，key不存在，将key设置为default值，并返回default这个值。
  \end{itemize}

\end{mdframed}

\subsection{第10题}
以下哪个是Python中用于科学计算与可视化的第三方库（）

\begin{enumerate}[label=\Alph*.]
  \item jieba
  \item scipy
  \item request
  \item random
\end{enumerate}

\subsection{第11题}
以下选项中，不是Python对文件的打开模式的是（）

\begin{enumerate}[label=\Alph*.]
  \item 'r'
  \item 'c'
  \item 'w'
  \item '+'
\end{enumerate}

\begin{mdframed}[linewidth=1pt, linecolor=black]

  \textbf{\color{red}【笔记】}

  正确答案：B

  Python文件打开模式说明：
  \begin{itemize}
    \item 'r'：只读模式（默认），文件必须存在
    \item 'w'：写入模式，会覆盖原文件内容，如果文件不存在则创建
    \item '+'：可读写模式，必须与其他模式组合使用，如'r+'、'w+'
    \item 'c'：不是Python的文件打开模式，此选项为错误选项
  \end{itemize}

  其他常见的文件打开模式：
  \begin{itemize}
    \item 'a'：追加模式，在文件末尾写入，如果文件不存在则创建
    \item 'x'：独占创建模式，文件必须不存在
    \item 'b'：二进制模式，与其他模式组合使用，如'rb'、'wb'
    \item 't'：文本模式（默认），与其他模式组合使用，如'rt'、'wt'
  \end{itemize}

\end{mdframed}

\subsection{第12题}
关于下面代码中的变量x，以下选项中描述正确的是（）
\begin{lstlisting}
fo = open(fname, "r")
for x in fo:
    print(x)
fo.close()
\end{lstlisting}

\begin{enumerate}[label=\Alph*.]
  \item 变量x表示文件中的一组字符
  \item 变量x表示文件中的一行字符
  \item 变量x表示文件中的一个字符
  \item 变量x表示文件中的全体字符
\end{enumerate}

\begin{mdframed}[linewidth=1pt, linecolor=black]

  \textbf{\color{red}【笔记】}

  正确答案：B

  代码分析：
  \begin{itemize}
    \item open(fname, "r")以只读模式打开文件，返回文件对象fo
    \item 在Python中，当使用for循环遍历文件对象时，每次迭代获得的是文件的一行内容
    \item 变量x在每次循环中存储文件的一行字符（包括行尾的换行符）
    \item print(x)会输出文件的每一行内容
    \item fo.close()关闭文件
  \end{itemize}

  知识点：
  \begin{itemize}
    \item 文件对象是可迭代的，for循环遍历文件对象时逐行读取
    \item 每行内容包含该行的所有字符以及行尾的换行符
    \item 这种方式适合处理大文件，因为每次只读取一行到内存中
  \end{itemize}

\end{mdframed}

\subsection{第13题}
以上代码输出结果为（）
\begin{lstlisting}
for i in "Python":
    print(i, end=", ")
\end{lstlisting}

\begin{enumerate}[label=\Alph*.]
  \item P*y*t*h*o*n*
  \item P,y,t,h,o,n,
  \item P y t h o n
  \item Python
\end{enumerate}

\begin{mdframed}[linewidth=1pt, linecolor=black]

  \textbf{\color{red}【笔记】}

  正确答案：B

  代码分析：
  \begin{itemize}
    \item for循环遍历字符串"Python"，每次取出一个字符赋值给变量i
    \item print(i, end=", ")中的end参数指定输出结束时不换行，而是输出", "
    \item 因此每个字符后面都会跟着逗号和空格
    \item 最终输出结果为：P, y, t, h, o, n,
  \end{itemize}

  知识点：print函数的end参数用于指定输出结束时的字符，默认为换行符\texttt{\\n}。

\end{mdframed}

\subsection{第14题}
执行Python语句\texttt{nums=set([1,2,2,3,3,3,4])}和\texttt{print(len(nums))}的结果是（）

\begin{enumerate}[label=\Alph*.]
  \item 1
  \item 2
  \item 4
  \item 7
\end{enumerate}

\begin{mdframed}[linewidth=1pt, linecolor=black]

  \textbf{\color{red}【笔记】}

  正确答案：C

  代码分析：
  \begin{itemize}
    \item 列表[1,2,2,3,3,3,4]包含7个元素，其中有重复元素
    \item set()函数将列表转换为集合，集合具有元素唯一性，会自动去除重复元素
    \item 转换后的集合nums包含元素\{1,2,3,4\}
    \item len(nums)返回集合中元素的个数，即4
  \end{itemize}

  知识点：
  \begin{itemize}
    \item set是Python的内置数据类型，表示无序且不重复的元素集合
    \item set会自动去除重复元素，保持元素的唯一性
    \item len()函数返回序列或集合中元素的个数
  \end{itemize}

\end{mdframed}

\subsection{第15题}
如下代码的输出为（）
\begin{lstlisting}
import re
s = 'a bc abc abbb abbbbbca'
re.findall('ab*', s)
\end{lstlisting}

\begin{enumerate}[label=\Alph*.]
  \item \verb|[|'ab', 'ab', 'ab'\verb|]|
  \item \verb|[|'ab', 'abbb', 'abbbb'\verb|]|
  \item \verb|[|'a', 'ab', 'abbb', 'abbbb', 'a'\verb|]|
  \item \verb|[|'a', 'ab', 'abbb', 'abbbb'\verb|]|
\end{enumerate}

\begin{mdframed}[linewidth=1pt, linecolor=black]

  \textbf{\color{red}【笔记】}

  正确答案：C

  \textbf{正则表达式详解：}

  \textbf{1. 基本概念：}
  \begin{itemize}
    \item 正则表达式（Regular Expression）是用于匹配字符串的强大工具
    \item 'ab*'表示匹配字母'a'后跟0个或多个字母'b'
    \item re.findall()返回字符串中所有匹配模式的子字符串列表
  \end{itemize}

  \textbf{2. 量词详解：}
  \begin{itemize}
    \item *：匹配前面字符0次或多次（贪婪匹配）
    \item +：匹配前面字符1次或多次
    \item ?：匹配前面字符0次或1次
    \item \{n\}：恰好匹配n次
    \item \{n,\}：至少匹配n次
    \item \{n,m\}：匹配n到m次
  \end{itemize}

  \textbf{3. 匹配过程分析：}
  \begin{lstlisting}
字符串：'a bc abc abbb abbbbbca'
模式：'ab*'

匹配过程：
位置0: 'a' -> 匹配'a'（0个b）✓
位置2: 'b' -> 不匹配（没有前导'a'）
位置3: 'c' -> 不匹配
位置5: 'a' -> 开始匹配
位置6: 'b' -> 匹配'ab'（1个b）✓
位置7: 'c' -> 结束匹配
位置9: 'a' -> 开始匹配
位置10-12: 'bbb' -> 匹配'abbb'（3个b）✓
位置14: 'a' -> 开始匹配
位置15-18: 'bbbb' -> 匹配'abbbb'（4个b）✓
位置19: 'b' -> 继续匹配但遇到'c'
位置20: 'c' -> 结束匹配
位置21: 'a' -> 匹配'a'（0个b）✓

结果：['a', 'ab', 'abbb', 'abbbb', 'a']
  \end{lstlisting}

  \textbf{4. 正则表达式语法详解：}

  \textbf{4.1 基本字符类：}
  \begin{lstlisting}
.     # 匹配任意字符（除换行符）
\d    # 匹配数字[0-9]
\D    # 匹配非数字[^0-9]
\w    # 匹配字母、数字、下划线[a-zA-Z0-9_]
\W    # 匹配非字母数字下划线[^a-zA-Z0-9_]
\s    # 匹配空白字符（空格、制表符、换行符等）
\S    # 匹配非空白字符
  \end{lstlisting}

  \textbf{4.2 自定义字符类：}
  \begin{lstlisting}
[abc]     # 匹配a、b或c中的任意一个
[a-z]     # 匹配任意小写字母
[A-Z]     # 匹配任意大写字母
[0-9]     # 匹配任意数字（等同于\d）
[a-zA-Z]  # 匹配任意字母
[^abc]    # 匹配除a、b、c外的任意字符（取反）
[a-z0-9]  # 匹配小写字母或数字
  \end{lstlisting}

  \textbf{4.3 量词（重复次数）：}
  \begin{lstlisting}
*         # 匹配前面的字符0次或多次
+         # 匹配前面的字符1次或多次
?         # 匹配前面的字符0次或1次
{n}       # 匹配前面的字符恰好n次
{n,}      # 匹配前面的字符至少n次
{n,m}     # 匹配前面的字符n到m次
*?        # 非贪婪匹配0次或多次
+?        # 非贪婪匹配1次或多次
??        # 非贪婪匹配0次或1次
  \end{lstlisting}

  \textbf{4.4 锚点（位置匹配）：}
  \begin{lstlisting}
^         # 匹配字符串开头
$         # 匹配字符串结尾
\b        # 匹配单词边界
\B        # 匹配非单词边界
\A        # 匹配字符串绝对开头
\Z        # 匹配字符串绝对结尾
  \end{lstlisting}

  \textbf{4.5 分组和选择：}
  \begin{lstlisting}
()        # 分组，创建捕获组
(?:...)   # 非捕获组
|         # 或者（选择）
\1, \2    # 反向引用第1、2个捕获组
(?P<name>...)  # 命名捕获组
(?P=name)      # 引用命名捕获组
  \end{lstlisting}

  \textbf{4.6 特殊字符转义：}
  \begin{lstlisting}
\.        # 匹配字面意思的点号
\*        # 匹配字面意思的星号
\+        # 匹配字面意思的加号
\?        # 匹配字面意思的问号
\[        # 匹配字面意思的左方括号
\]        # 匹配字面意思的右方括号
\\        # 匹配字面意思的反斜杠
\^        # 匹配字面意思的插入符
\$        # 匹配字面意思的美元符号
  \end{lstlisting}

  \textbf{4.7 实用示例：}
  \begin{lstlisting}
# 匹配邮箱
r'\w+@\w+\.\w+'

# 匹配手机号（简单版）
r'1[3-9]\d{9}'

# 匹配IP地址
r'\d{1,3}\.\d{1,3}\.\d{1,3}\.\d{1,3}'

# 匹配HTML标签
r'<[^>]+>'

# 匹配中文字符
r'[\u4e00-\u9fa5]+'

# 匹配URL
r'https?://[^\s]+'
  \end{lstlisting}

  \textbf{5. re模块常用函数：}
  \begin{lstlisting}
import re

text = 'Python 3.8 is great!'

# 查找所有匹配
re.findall(r'\d+', text)        # ['3', '8']

# 搜索第一个匹配
match = re.search(r'\d+', text)  # 返回Match对象或None

# 从字符串开头匹配
match = re.match(r'Python', text)  # 匹配成功

# 替换
result = re.sub(r'\d+', 'X', text)  # 'Python X.X is great!'

# 分割
parts = re.split(r'\s+', text)   # ['Python', '3.8', 'is', 'great!']
  \end{lstlisting}

  \textbf{6. 贪婪匹配vs非贪婪匹配：}
  \begin{lstlisting}
text = 'abbbbbb'

# 贪婪匹配（默认）
re.findall('ab*', text)      # ['abbbbbb'] - 匹配尽可能多的b

# 非贪婪匹配
re.findall('ab*?', text)     # ['a'] - 匹配尽可能少的b
  \end{lstlisting}

  \textbf{知识点总结：}
  \begin{itemize}
    \item 理解量词的含义和用法
    \item 掌握正则表达式的匹配原理
    \item 熟悉re模块的常用函数
    \item 区分贪婪匹配和非贪婪匹配
  \end{itemize}

\end{mdframed}

\subsection{第16题}
以下不能创建一个元组的语句是（）

\begin{enumerate}[label=\Alph*.]
  \item tup1 = ()
  \item tup2 = 1,
  \item tup3 = (1)
  \item dict4 = tuple("123")
\end{enumerate}

\begin{mdframed}[linewidth=1pt, linecolor=black]

  \textbf{\color{red}【笔记】}

  正确答案：C

  各选项分析：
  \begin{itemize}
    \item A选项：tup1 = () 创建一个空元组，正确
    \item B选项：tup2 = 1, 创建包含一个元素的元组，逗号是关键，正确
    \item C选项：tup3 = (1) 仅仅是给数字1加括号，结果是整数1，不是元组
    \item D选项：tuple("123") 将字符串转换为元组('1', '2', '3')，正确
  \end{itemize}

  知识点：
  \begin{itemize}
    \item 创建元组时，逗号是关键标识符，不是括号
    \item 单个元素的元组必须在元素后加逗号：(1,) 或 1,
    \item 空元组可以用 () 或 tuple() 创建
    \item tuple()函数可以将其他可迭代对象转换为元组
  \end{itemize}

\end{mdframed}

\subsection{第17题}
在PyCharm中运行整个程序的默认快捷键是（）

\begin{enumerate}[label=\Alph*.]
  \item Shift+F10
  \item Ctrl+F10
  \item Shift+Ctrl+F10
  \item Shift+Ctrl+Enter
\end{enumerate}

\begin{mdframed}[linewidth=1pt, linecolor=black]

  \textbf{\color{red}【笔记】}

  正确答案：A

  PyCharm常用快捷键：
  \begin{itemize}
    \item Shift+F10：运行整个程序（Run）
    \item Ctrl+F10：运行当前文件
    \item Shift+Ctrl+F10：运行当前光标所在的配置
    \item Shift+Ctrl+Enter：完成当前语句
  \end{itemize}

  其他常用快捷键：
  \begin{itemize}
    \item Shift+F9：调试程序（Debug）
    \item Ctrl+Shift+F9：调试当前文件
    \item F8：单步执行（调试时）
    \item F9：继续执行（调试时）
  \end{itemize}

\end{mdframed}

\subsection{第18题}
下面的语句中（）用来把路径path设置为默认路径

\begin{enumerate}[label=\Alph*.]
  \item os.chdir(path)
  \item os.mkdir(path)
  \item os.isdir(path)
  \item os.listdir(path)
\end{enumerate}

\begin{mdframed}[linewidth=1pt, linecolor=black]

  \textbf{\color{red}【笔记】}

  正确答案：A

  os模块常用函数说明：
  \begin{itemize}
    \item os.chdir(path)：改变当前工作目录到指定路径，相当于cd命令
    \item os.mkdir(path)：创建一个新目录
    \item os.isdir(path)：判断路径是否为目录，返回True或False
    \item os.listdir(path)：返回指定目录下的文件和目录列表
  \end{itemize}

  其他相关函数：
  \begin{itemize}
    \item os.getcwd()：获取当前工作目录
    \item os.makedirs(path)：递归创建目录
    \item os.path.exists(path)：判断路径是否存在
    \item os.path.join()：连接路径
  \end{itemize}

\end{mdframed}

\subsection{第19题}
在Python中，下列说法正确的是（）

\begin{enumerate}[label=\Alph*.]
  \item 0xad是合法的十六进制数字表示形式
  \item 3+4j不是合法的Python表达式
  \item 可以使用if作为变量名
  \item 0o12f是合法的八进制数字
\end{enumerate}

\begin{mdframed}[linewidth=1pt, linecolor=black]

  \textbf{\color{red}【笔记】}

  正确答案：A

  各选项分析：
  \begin{itemize}
    \item A选项：0xad是合法的十六进制数字，0x前缀表示十六进制，ad是有效的hex数字
    \item B选项：3+4j是合法的复数表达式，j表示虚数单位
    \item C选项：if是Python保留字（关键字），不能用作变量名
    \item D选项：0o12f不合法，八进制数字只能包含0-7，f不是有效的八进制数字
  \end{itemize}

  Python数字表示形式：
  \begin{itemize}
    \item 十进制：直接写数字，如123
    \item 二进制：0b前缀，如0b1010
    \item 八进制：0o前缀，如0o777（只能包含0-7）
    \item 十六进制：0x前缀，如0xff（可包含0-9和a-f）
    \item 复数：实部+虚部j，如3+4j
  \end{itemize}

\end{mdframed}

\section{判断题}

\subsection{第1题}
定义函数时，即便该函数不需要接收任何参数，也必须保留一对空的圆括号来表示这是一个函数。（\quad）

\begin{mdframed}[linewidth=1pt, linecolor=black]
  \textbf{\color{red}【笔记】}

  正确答案：正确 (T)

  解释：
  \begin{itemize}
    \item Python函数定义语法：\texttt{def function\_name():}
    \item 即使函数不需要参数，圆括号 \texttt{()} 也是必须的
    \item 圆括号是函数定义的语法要求，用于区分函数和变量
    \item 示例：\texttt{def hello(): print("Hello")}
  \end{itemize}

\end{mdframed}

\subsection{第2题}
已知\texttt{ls=[2, 4, 6]}，那么执行语句\texttt{ls.append(8)}之后，ls的内存地址不变。（\quad）

\begin{mdframed}[linewidth=1pt, linecolor=black]
  \textbf{\color{red}【笔记】}

  正确答案：正确 (T)

  解释：
  \begin{itemize}
    \item \texttt{append()}方法是原地修改操作（in-place operation）
    \item 列表对象本身不会被重新创建，只是在原有内存空间中添加新元素
    \item 可以用\texttt{id()}函数验证：\texttt{id(ls)}在append前后相同
    \item 类似的原地修改方法还有：\texttt{extend()}, \texttt{insert()}, \texttt{remove()}等
  \end{itemize}

  对比：重新赋值操作如\texttt{ls = ls + [8]}会创建新的列表对象，内存地址会改变。

\end{mdframed}

\subsection{第3题}
continue语句在一旦在循环结构里被执行，将使得当前循环的整个循环提前结束。（\quad）

\begin{mdframed}[linewidth=1pt, linecolor=black]
  \textbf{\color{red}【笔记】}

  正确答案：错误 (F)

  解释：
  \begin{itemize}
    \item \texttt{continue}语句只是跳过当前循环迭代，继续下一次迭代
    \item \texttt{continue}不会终止整个循环，循环仍会继续执行
    \item \texttt{break}语句才会提前终止整个循环
  \end{itemize}

  示例对比：
  \begin{itemize}
    \item 使用\texttt{continue}：跳过满足条件的迭代，继续其他迭代
    \item 使用\texttt{break}：一旦满足条件就立即退出整个循环
  \end{itemize}

\end{mdframed}

\subsection{第4题}
\texttt{a,b=b,a}可以实现a和b值的交换。（\quad）

\begin{mdframed}[linewidth=1pt, linecolor=black]
  \textbf{\color{red}【笔记】}

  正确答案：正确 (T)

  解释：
  \begin{itemize}
    \item Python支持多重赋值（multiple assignment）
    \item 右侧\texttt{b,a}先被打包成元组\texttt{(b,a)}
    \item 然后元组被解包赋值给左侧的\texttt{a,b}
    \item 这种方式简洁高效，无需临时变量
  \end{itemize}

  传统交换方式对比：
  \begin{itemize}
    \item 传统方式：\texttt{temp=a; a=b; b=temp}（需要临时变量）
    \item Python方式：\texttt{a,b=b,a}（一行代码完成）
  \end{itemize}

\end{mdframed}

\subsection{第5题}
函数是封装了一些独立的功能，可以直接调用，Python内置了许多函数，同时可以自建函数来使用。（\quad）

\begin{mdframed}[linewidth=1pt, linecolor=black]
  \textbf{\color{red}【笔记】}

  正确答案：正确 (T)

  解释：
  \begin{itemize}
    \item 函数是代码重用和模块化的基础
    \item 函数封装特定功能，提高代码可读性和可维护性
    \item Python内置函数如：\texttt{print()}, \texttt{len()}, \texttt{max()}, \texttt{min()}等
    \item 用户可以使用\texttt{def}关键字自定义函数
  \end{itemize}

  函数的优点：
  \begin{itemize}
    \item 代码重用：避免重复编写相同代码
    \item 模块化：将复杂问题分解为小问题
    \item 易于测试和调试
  \end{itemize}

\end{mdframed}

\subsection{第6题}
使用内置函数\texttt{open()}且以'w'模式打开的文件，文件指针默认指向文件尾。（\quad）

\begin{mdframed}[linewidth=1pt, linecolor=black]
  \textbf{\color{red}【笔记】}

  正确答案：错误 (F)

  解释：
  \begin{itemize}
    \item 'w'模式（写入模式）文件指针指向文件开头
    \item 'w'模式会清空原文件内容，从头开始写入
    \item 'a'模式（追加模式）文件指针才指向文件尾
  \end{itemize}

  文件打开模式对比：
  \begin{itemize}
    \item 'r'：只读模式，指针在文件开头
    \item 'w'：写入模式，清空文件，指针在文件开头
    \item 'a'：追加模式，指针在文件末尾
    \item 'r+'：读写模式，指针在文件开头
  \end{itemize}

\end{mdframed}

\subsection{第7题}
利用\texttt{print()}格式化输出，\verb|{2:f}|能够控制浮点数的小数点后保留两位。（\quad）

\begin{mdframed}[linewidth=1pt, linecolor=black]
  \textbf{\color{red}【笔记】}

  正确答案：错误 (F)

  解释：
  \begin{itemize}
    \item \verb|{2:f}|不能控制小数点后保留两位
    \item 2表示参数索引（第3个参数，从0开始计数）
    \item f表示浮点数格式，但没有指定精度
    \item 要保留两位小数需要使用\verb|{2:.2f}|
  \end{itemize}

  正确的格式化语法：
  \begin{itemize}
    \item \verb|{2:.2f}|：第3个参数保留2位小数
    \item \verb|{:.2f}|：当前位置参数保留2位小数
    \item \verb|{2:f}|：第3个参数浮点数格式，但精度由系统默认决定
  \end{itemize}

  \label{string:format}
  \textbf{Python print()格式化输出详解：}

  \textbf{1. 基本格式化语法}
  \begin{itemize}
    \item 基本形式：\verb|print("格式字符串".format(参数))|
    \item f-string形式：\verb|print(f"格式字符串{变量}")| （Python 3.6+）
    \item \%格式化：\verb|print("格式字符串" % 参数)| （较老的方式）
    \end{itemize}

    \textbf{2. 格式说明符详解}
    \begin{itemize}
      \item \verb|{索引:格式说明符}|
      \item 索引：0, 1, 2... 或省略（按顺序）
      \item 格式说明符：[填充字符][对齐][宽度][.精度][类型]
    \end{itemize}

    \textbf{3. 数字格式化类型}
    \begin{itemize}
      \item \texttt{d}：十进制整数
      \item \texttt{f}：浮点数（默认6位小数）
      \item \texttt{.nf}：浮点数保留n位小数
      \item \texttt{e}：科学计数法
      \item \texttt{g}：自动选择f或e格式
      \item \texttt{o}：八进制，\texttt{x}：十六进制
    \end{itemize}

    \textbf{4. 字符串格式化}
    \begin{itemize}
      \item \texttt{s}：字符串格式
      \item \verb|{:10s}|：字符串宽度10，左对齐
      \item \verb|{:>10s}|：右对齐，\verb|{:^10s}|：居中对齐
      \item \verb|{:*^10s}|：用*填充的居中对齐
    \end{itemize}

    \textbf{5. 格式化示例}
    \begin{itemize}
      \item \verb|print("{:.2f}".format(3.14159))| → 3.14
      \item \verb|print("{:10.2f}".format(3.14))| → \quad\quad\quad 3.14
      \item \verb|print("{:0>5d}".format(42))| → 00042
      \item \verb|print(f"{name:^10s}")| → 变量name居中显示
    \end{itemize}

  \end{mdframed}

  \subsection{第8题}
  定义类时实现了\texttt{\_\_pow\_\_()}方法，该类对象即可支持运算符**。（\quad）

  \begin{mdframed}[linewidth=1pt, linecolor=black]
    \textbf{\color{red}【笔记】}

    正确答案：正确 (T)

    解释：
    \begin{itemize}
      \item \texttt{\_\_pow\_\_()}是Python的魔法方法（magic method）
      \item 实现该方法后，类对象可以使用**运算符进行幂运算
      \item 调用\texttt{obj ** n}等价于调用\texttt{obj.\_\_pow\_\_(n)}
    \end{itemize}

    其他常用魔法方法：
    \begin{itemize}
      \item \texttt{\_\_add\_\_()}：支持+运算符
      \item \texttt{\_\_sub\_\_()}：支持-运算符
      \item \texttt{\_\_mul\_\_()}：支持*运算符
      \item \texttt{\_\_str\_\_()}：支持字符串转换
    \end{itemize}

  \end{mdframed}

  \subsection{第9题}
  通过对象和类名都可以调用类方法和静态方法。（\quad）

  \begin{mdframed}[linewidth=1pt, linecolor=black]
    \textbf{\color{red}【笔记】}

    正确答案：正确 (T)

    解释：
    \begin{itemize}
      \item 类方法（\texttt{@classmethod}）：可通过类名和对象调用
      \item 静态方法（\texttt{@staticmethod}）：可通过类名和对象调用
      \item 这两种方法不依赖于特定的实例状态
    \end{itemize}

    方法调用方式对比：
    \begin{itemize}
      \item 实例方法：只能通过对象调用，需要\texttt{self}参数
      \item 类方法：类名.方法() 或 对象.方法()，接收\texttt{cls}参数
      \item 静态方法：类名.方法() 或 对象.方法()，无特殊参数要求
    \end{itemize}

    推荐做法：类方法和静态方法优先使用类名调用，语义更清晰。

    \textbf{类方法与静态方法详细区别：}

    \textbf{1. 定义方式}
    \begin{itemize}
      \item 类方法：使用\texttt{@classmethod}装饰器
      \item 静态方法：使用\texttt{@staticmethod}装饰器
    \end{itemize}

    \textbf{2. 参数区别}
    \begin{itemize}
      \item 类方法：第一个参数必须是\texttt{cls}（代表类本身）
      \item 静态方法：没有特殊的第一个参数要求
      \item 实例方法：第一个参数必须是\texttt{self}（代表实例本身）
    \end{itemize}

    \textbf{3. 访问权限}
    \begin{itemize}
      \item 类方法：可以访问类属性，不能直接访问实例属性
      \item 静态方法：不能访问类属性和实例属性
      \item 实例方法：可以访问类属性和实例属性
    \end{itemize}

    \textbf{4. 使用场景}
    \begin{itemize}
      \item 类方法：替代构造函数、操作类属性、工厂方法
      \item 静态方法：与类相关但不需要访问类/实例数据的工具函数
      \item 实例方法：操作具体实例的数据和行为
    \end{itemize}

    \textbf{5. 代码示例}
    \begin{lstlisting}[linewidth=\textwidth, breaklines=true, numbers=none]
class MyClass:
    class_var = "类属性"

    def __init__(self, value):
        self.instance_var = value

    @classmethod
    def class_method(cls):
        return f"访问: {cls.class_var}"

    @staticmethod
    def static_method():
        return "静态方法"

    def instance_method(self):
        return f"实例: {self.instance_var}"

# 调用方式
MyClass.class_method()    # 类方法
MyClass.static_method()   # 静态方法
obj = MyClass("值")
obj.class_method()        # 对象调用类方法
obj.static_method()       # 对象调用静态方法
    \end{lstlisting}
  \end{mdframed}

\section{填空题}

\subsection{第1题}
有如下Python程序：
\begin{lstlisting}
def f(a, b):
    if b == 0:
        print(a)
    else:
        f(b, a % b)
\end{lstlisting}
则\texttt{print(f(9,6))}的输出结果是\underline{\hspace{2cm}}、\underline{\hspace{2cm}}

\begin{mdframed}[linewidth=1pt, linecolor=black]

  \textbf{\color{red}【笔记】}

  正确答案：\underline{3}, \underline{None}

  解释：

  \qquad 这是一个经典的辗转相除法求最大公约数（GCD）的递归代码。
  当b为0时，a就等于gcd(a,b)的值。由于没有返回值，所以print(f(a,b))输出None。

\end{mdframed}

\subsection{第2题}
表达式\texttt{sorted(['abc', 'acd', 'ade'], key=lambda x:(x[0],x[2]))}的值为\\
\underline{\hspace{8cm}}

\begin{mdframed}[linewidth=1pt, linecolor=black]

  \textbf{\color{red}【笔记】}

  正确答案：\underline{['abc', 'acd', 'ade']}

  解释：

  \qquad sorted函数使用key参数来指定排序的依据。这里的key是lambda x:(x[0],x[2])，
  意味着对每个字符串取第0个字符和第2个字符组成元组作为排序键。

  \qquad 分析每个字符串的排序键：
  \begin{itemize}
    \item 'abc': (x[0], x[2]) = ('a', 'c')
    \item 'acd': (x[0], x[2]) = ('a', 'd')
    \item 'ade': (x[0], x[2]) = ('a', 'e')
  \end{itemize}

  \qquad 由于第0个字符都是'a'，按第2个字符排序：'c' < 'd' < 'e'，
  所以原顺序已经是正确的排序结果。

\end{mdframed}

\subsection{第3题}
表达式\verb|'The first:{1:.3f}, the second is {0:b}'.format(2,3.1415926)|的值为\\
\underline{\hspace{8cm}}

\begin{mdframed}[linewidth=1pt, linecolor=black]

  \textbf{\color{red}【笔记】}

  正确答案：\underline{'The first:3.142, the second is 10'}

  解释：

  \qquad 这是字符串的format方法格式化，format(2, 3.1415926)传入两个参数：
  \begin{itemize}
    \item 参数0：2
    \item 参数1：3.1415926
  \end{itemize}

  \qquad 格式化说明：
  \begin{itemize}
    \item \{1:.3f\}：取第1个参数3.1415926，格式化为保留3位小数的浮点数 → 3.142
    \item \{0:b\}：取第0个参数2，格式化为二进制 → 10
  \end{itemize}

  \qquad 因此最终结果为：'The first:3.142, the second is 10'

  详细字符串格式化解释在\ref{string:format}
\end{mdframed}

\subsection{第4题}
假设正则表达式模块re已导入，那么表达式\texttt{re.sub('\textbackslash d+', '1', 'a12345bbbbb67c890d0e')}的值为\\
\underline{\hspace{8cm}}

\begin{mdframed}[linewidth=1pt, linecolor=black]

  \textbf{\color{red}【笔记】}

  正确答案：\underline{'a1bbbbb1c1d1e'}

  解释：

  \qquad re.sub(pattern, replacement, string)函数用于字符串替换：
  \begin{itemize}
    \item pattern: '\textbackslash d+' 即 \textbackslash d+，匹配一个或多个连续数字
    \item replacement: '1'，替换成的字符串
    \item string: 'a12345bbbbb67c890d0e'，要处理的原字符串
  \end{itemize}

  \qquad 匹配过程：
  \begin{itemize}
    \item 原字符串：'a12345bbbbb67c890d0e'
    \item \textbackslash d+ 匹配到的数字序列：12345, 67, 890, 0
    \item 每个数字序列都被替换为'1'
  \end{itemize}

  \qquad 替换结果：'a1bbbbb1c1d1e'

  \textbf{re模块的常见函数及用法}

  \begin{itemize}
    \item \textbf{re.match(pattern, string[, flags])}：从字符串开头开始匹配
      \begin{itemize}
        \item 只匹配字符串开头，如果开头不符合则返回None
        \item 示例：\texttt{re.match(r'\textbackslash d+', '123abc')} → 匹配到'123'
        \item 示例：\texttt{re.match(r'\textbackslash d+', 'abc123')} → 返回None
      \end{itemize}

    \item \textbf{re.search(pattern, string[, flags])}：在整个字符串中搜索第一个匹配
      \begin{itemize}
        \item 扫描整个字符串，返回第一个匹配的结果
        \item 示例：\texttt{re.search(r'\textbackslash d+', 'abc123def456')} → 匹配到'123'
      \end{itemize}

    \item \textbf{re.findall(pattern, string[, flags])}：查找所有匹配项
      \begin{itemize}
        \item 返回一个列表，包含所有匹配的字符串
        \item 示例：\texttt{re.findall(r'\textbackslash d+', 'a12b34c56')} → ['12', '34', '56']
      \end{itemize}

    \item \textbf{re.finditer(pattern, string[, flags])}：返回匹配对象的迭代器
      \begin{itemize}
        \item 返回一个迭代器，每个元素是一个Match对象
        \item 适用于大文本，节省内存
        \item 示例：\texttt{for m in re.finditer(r'\textbackslash d+', 'a12b34'): print(m.group())}
      \end{itemize}

    \item \textbf{re.sub(pattern, repl, string[, count=0, flags=0])}：替换匹配项
      \begin{itemize}
        \item 将匹配的部分替换为指定字符串
        \item count参数指定最多替换次数，0表示全部替换
        \item 示例：\texttt{re.sub(r'\textbackslash d+', 'X', 'a12b34')} → 'aXbX'
      \end{itemize}

    \item \textbf{re.split(pattern, string[, maxsplit=0, flags=0])}：按模式分割字符串
      \begin{itemize}
        \item 使用正则表达式作为分隔符分割字符串
        \item 示例：\texttt{re.split(r'[,;]', 'a,b;c,d')} → ['a', 'b', 'c', 'd']
      \end{itemize}

    \item \textbf{re.compile(pattern[, flags])}：编译正则表达式
      \begin{itemize}
        \item 将正则表达式编译成Pattern对象，提高重复使用的效率
        \item 示例：\texttt{p = re.compile(r'\textbackslash d+'); p.findall('a12b34')} → ['12', '34']
      \end{itemize}
  \end{itemize}

  \textbf{常用的flags参数：}
  \begin{itemize}
    \item \texttt{re.I} 或 \texttt{re.IGNORECASE}：忽略大小写
    \item \texttt{re.M} 或 \texttt{re.MULTILINE}：多行模式，\^{}和\$匹配每行的开头和结尾
    \item \texttt{re.S} 或 \texttt{re.DOTALL}：使.匹配包括换行符在内的所有字符
    \item \texttt{re.X} 或 \texttt{re.VERBOSE}：忽略空白和注释，可以写更清晰的正则表达式
  \end{itemize}

\end{mdframed}

\subsection{第5题}
下面代码的功能是，随机生成20个介于[1,50]之间的整数，然后统计每个整数出现频率。请把缺少的代码补全。
\begin{lstlisting}
import random
x = [random.______(1, 50) for i in range(20)]
d = dict()
for i in x:
    d[i] = d.get(i, ______) + 1
for k, v in d.items():
    print(k, v)
\end{lstlisting}
上述横线处分别填写\underline{\hspace{3cm}}、\underline{\hspace{3cm}}

\begin{mdframed}[linewidth=1pt, linecolor=black]

  \textbf{\color{red}【笔记】}

  正确答案：\underline{randint}、\underline{0}

  解释：

  \qquad 第一个空：random.\underline{randint}(1, 50)
  \begin{itemize}
    \item randint(a, b)函数生成[a, b]范围内的随机整数（包含端点）
    \item 这里需要生成[1, 50]之间的整数，所以用randint
  \end{itemize}

  \qquad 第二个空：d.get(i, \underline{0})
  \begin{itemize}
    \item dict.get(key, default)方法：如果key存在返回对应值，否则返回default
    \item 统计频率时，如果数字第一次出现，频率应该是0+1=1
    \item 如果数字已经存在，频率就是原来的值+1
  \end{itemize}

  \qquad 代码逻辑：生成20个随机数→遍历每个数→统计每个数的出现次数→输出结果

\end{mdframed}

\section{简答题}

\subsection{第1题}
在Python的运算符中，|属于何种类型运算符？在集合set中|有什么作用？并分析逻辑运算符"or"的短路求值特性。

\begin{mdframed}[linewidth=1pt, linecolor=black]

  \textbf{\color{red}【笔记】}

  \textbf{答案：}

  \textbf{1. |运算符的类型：}
  \qquad |属于\textbf{位运算符}（按位或运算符）。它对整数的二进制位进行按位或运算。

  \textbf{2. 在集合set中|的作用：}
  \qquad 在集合中，|表示\textbf{并集运算}。例如：
  \begin{lstlisting}
set1 = {1, 2, 3}
set2 = {3, 4, 5}
result = set1 | set2  # 结果：{1, 2, 3, 4, 5}
  \end{lstlisting}

  \textbf{3. or的短路求值特性：}
  \qquad \textbf{短路求值}是指当第一个操作数能够确定整个表达式的结果时，就不再计算第二个操作数。

  \qquad 对于or运算符：
  \begin{itemize}
    \item 如果第一个操作数为True，则不评估第二个操作数，直接返回第一个操作数的值
    \item 如果第一个操作数为False，则返回第二个操作数的值
  \end{itemize}

  \qquad 示例：
  \begin{lstlisting}
def func():
    print("函数被调用了")
    return True

result = True or func()  # func()不会被调用，不会打印
  \end{lstlisting}

\end{mdframed}

\subsection{第2题}
\texttt{zip()}函数可以将两个或多个可迭代类型(元组、列表等)组合为一个关系对(元组)，并返回包含这些元素的zip对象，该函数非常适合生成键值对。

请根据以下示例数据：
\begin{itemize}
  \item \verb|[|(1, 'A'), (2, 'B'), (3, 'C'), (4, 'D')\verb|]|
  \item \verb|{1: 'A', 2: 'B', 3: 'C', 4: 'D'}|(1分)
  \item \verb|[|'A', 'B', 'C', 'D'\verb|]|
\end{itemize}

\subsection{第3题}
请列举pip常用的5个命令及其作用。


\section{代码阅读题}

\subsection{第1题}
阅读下面的代码，并分析假设文件"E:$\backslash$writeTest.txt"不存在的情况下两段代码可能发生的问题。

\textbf{代码1：}
\begin{lstlisting}
fp = open(r'E:\writeTest.txt')
fp.write("python")
fp.close()
\end{lstlisting}

\textbf{代码2：}
\begin{lstlisting}
fp = open(r'E:\writeTest.txt', 'a+')
fp.write("python")
fp.close()
\end{lstlisting}

\begin{mdframed}[linewidth=1pt, linecolor=black]

  \textbf{\color{red}【笔记】}

  \textbf{问题分析：}

  \textbf{代码1存在的问题：}
  \begin{itemize}
    \item \texttt{open(r'E:\\writeTest.txt')}没有指定打开模式，默认为'r'（只读模式）
    \item 如果文件不存在，会抛出\texttt{FileNotFoundError}异常
    \item 即使文件存在，只读模式下调用\texttt{write()}方法也会抛出\texttt{io.UnsupportedOperation}异常
    \item 这段代码无法正常执行
  \end{itemize}

  \textbf{代码2的情况：}
  \begin{itemize}
    \item 使用'a+'模式（追加+读写模式）
    \item 如果文件不存在，会自动创建该文件
    \item 可以正常写入内容"python"
    \item 这段代码可以正常执行
  \end{itemize}

  \textbf{文件打开模式详细说明：}

  \textbf{基本模式：}
  \begin{itemize}
    \item \textbf{'r'}：只读模式（默认模式）
      \begin{itemize}
        \item 文件不存在时抛出FileNotFoundError
        \item 只能读取，不能写入
        \item 文件指针位于文件开头
      \end{itemize}

    \item \textbf{'w'}：写入模式
      \begin{itemize}
        \item 文件不存在时自动创建
        \item 文件存在时清空内容（覆盖）
        \item 只能写入，不能读取
        \item 文件指针位于文件开头
      \end{itemize}

    \item \textbf{'a'}：追加模式
      \begin{itemize}
        \item 文件不存在时自动创建
        \item 文件存在时保留原内容
        \item 只能写入，不能读取
        \item 文件指针位于文件末尾
      \end{itemize}
  \end{itemize}

  \textbf{复合模式（可读可写）：}
  \begin{itemize}
    \item \textbf{'r+'}：读写模式
      \begin{itemize}
        \item 文件必须存在，否则报错
        \item 可读可写，文件指针位于开头
        \item 写入会覆盖对应位置的内容
      \end{itemize}

    \item \textbf{'w+'}：写读模式
      \begin{itemize}
        \item 文件不存在时创建，存在时清空
        \item 可读可写，文件指针位于开头
      \end{itemize}

    \item \textbf{'a+'}：追加读写模式
      \begin{itemize}
        \item 文件不存在时创建，存在时保留内容
        \item 可读可写，写入时指针自动移到末尾
        \item 读取时可以移动指针到任意位置
      \end{itemize}
  \end{itemize}

  \textbf{二进制模式（在上述模式后加'b'）：}
  \begin{itemize}
    \item 'rb'、'wb'、'ab'、'r+b'、'w+b'、'a+b'
    \item 用于处理二进制文件（图片、视频、exe等）
    \item 返回bytes对象而不是字符串
  \end{itemize}

\end{mdframed}

\subsection{第2题}
代码分析

\textbf{(1)} 阅读下面代码，解释其功能并写出执行结果。
\begin{lstlisting}
def demo(*p):
    return sum(p)
print(demo(1, 2, 3, 4, 5))
print(demo(1, 2, 3))
\end{lstlisting}

\textbf{(2)} 阅读下面代码，解释其功能并写出执行结果。
\begin{lstlisting}
def Join(ls, sep=None):
    return (sep or ',').join(ls)
print(Join(['a', 'b', 'c']))
print(Join(['a', 'b', 'c'], ':'))
\end{lstlisting}

\begin{mdframed}[linewidth=1pt, linecolor=black]

  \textbf{\color{red}【笔记】}

  \textbf{(1) 代码分析：}

  \textbf{功能：}定义一个可变参数函数，计算所有传入参数的和。

  \textbf{代码解释：}
  \begin{itemize}
    \item \texttt{*p}表示可变参数，接收任意数量的位置参数，组成元组
    \item \texttt{sum(p)}计算元组中所有数字的和
    \item 函数返回计算结果
  \end{itemize}

  \textbf{执行结果：}
  \begin{lstlisting}
15
6
  \end{lstlisting}

  \textbf{(2) 代码分析：}

  \textbf{功能：}定义一个字符串连接函数，使用指定分隔符或默认逗号连接列表元素。

  \textbf{代码解释：}
  \begin{itemize}
    \item \texttt{sep=None}设置默认参数为None
    \item \texttt{(sep or ',')}利用短路求值：如果sep为None或False，则使用','
    \item \texttt{join(ls)}用分隔符连接列表中的字符串元素
  \end{itemize}

  \textbf{执行结果：}
  \begin{lstlisting}
a,b,c
a:b:c
  \end{lstlisting}

  \textbf{知识点总结：}
  \begin{itemize}
    \item 可变参数(\texttt{*args})的使用
    \item 默认参数和短路求值的应用
    \item 字符串\texttt{join()}方法的使用
  \end{itemize}

\end{mdframed}

\subsection{第3题}
给出如下代码：
\begin{lstlisting}
from random import randint
result = list()
while True:
    result.append(randint(1, 10))
    if len(result) == 20:
        break
print(result)
\end{lstlisting}

以上代码中，result为何种类型变量？程序是否能够正常执行，若不能，请解释原因；若能，请分析其执行结果。

\begin{mdframed}[linewidth=1pt, linecolor=black]

  \textbf{\color{red}【笔记】}

  \textbf{代码分析：}

  \textbf{1. result的类型：}
  \begin{itemize}
    \item \texttt{result = list()}创建了一个空列表
    \item 因此result是\textbf{列表(list)类型}变量
  \end{itemize}

  \textbf{2. 程序执行情况：}
  \begin{itemize}
    \item \textbf{程序可以正常执行}
    \item 代码语法正确，逻辑清晰
  \end{itemize}

  \textbf{3. 程序执行过程：}
  \begin{itemize}
    \item 创建空列表result
    \item 进入无限循环\texttt{while True}
    \item 每次循环生成一个[1,10]范围的随机整数并添加到列表
    \item 当列表长度达到20时，执行\texttt{break}跳出循环
    \item 打印包含20个随机数的列表
  \end{itemize}

  \textbf{4. 执行结果：}
  \begin{itemize}
    \item 输出：包含20个随机整数的列表
    \item 每个整数都在[1,10]范围内（包含1和10）
    \item 示例输出：[3, 7, 1, 9, 2, 5, 8, 4, 10, 6, 1, 3, 9, 7, 2, 8, 5, 4, 10, 6]
    \item 注意：由于是随机数，每次运行结果都不同
  \end{itemize}

  \textbf{知识点总结：}
  \begin{itemize}
    \item 列表的创建方法：\texttt{list()}和\texttt{[]}
    \item \texttt{append()}方法向列表末尾添加元素
    \item \texttt{len()}函数获取列表长度
    \item 无限循环\texttt{while True}配合\texttt{break}的使用
  \end{itemize}

\end{mdframed}

\subsection{第4题}
根据如下代码回答问题：
\begin{lstlisting}
class MyArray:
    def __init__(self, *args):
        if not args:
            self.__value = []
        else:
            for arg in args:
                self.__value = list(args)

    def __sub__(self, n):
        b = MyArray()
        b.__value = [item - n for item in self.__value]
        return b

    def __str__(self):
        return str(self.__value)

# 测试主程序
if __name__ == '__main__':
    m = MyArray(10, 3, 2, 5)
\end{lstlisting}

请回答以下问题：

\textbf{(1)} \texttt{\_\_init\_\_}中参数\texttt{*args}具有什么含义？

\textbf{(2)} 请说明\texttt{\_\_sub\_\_}方法的功能。并在测试主程序中调用该方法，使其输出为[6, -1, -2, 1]，写出调用代码。

\begin{mdframed}[linewidth=1pt, linecolor=black]

  \textbf{\color{red}【笔记】}

  \textbf{代码分析：}

  \textbf{(1) *args的含义：}
  \begin{itemize}
    \item \texttt{*args}是可变参数，表示接收任意数量的位置参数
    \item 在函数内部，args是一个元组，包含所有传入的参数
    \item 允许创建MyArray对象时传入0个或多个参数
    \item 例如：\texttt{MyArray()}、\texttt{MyArray(1)}、\texttt{MyArray(1,2,3)}都是合法的
  \end{itemize}

  \textbf{构造函数逻辑：}
  \begin{itemize}
    \item 如果没有传入参数(\texttt{not args})，创建空列表
    \item 如果传入参数，将args转换为列表赋值给\texttt{self.\_\_value}
    \item 注意：代码中的for循环是多余的，每次都会重新赋值
  \end{itemize}

  \textbf{(2) \_\_sub\_\_方法功能：}
  \begin{itemize}
    \item \texttt{\_\_sub\_\_}是Python的魔术方法，重载减法运算符'-'
    \item 功能：将数组中每个元素都减去指定的数n
    \item 返回一个新的MyArray对象，不修改原对象
    \item 实现：\texttt{[item - n for item in self.\_\_value]}列表推导式
  \end{itemize}

  \textbf{调用代码分析：}
  \begin{itemize}
    \item 当前：\texttt{m = MyArray(10, 3, 2, 5)}，即m.\_\_value = [10, 3, 2, 5]
    \item 要得到[6, -1, -2, 1]，需要每个元素减去4：
      \begin{itemize}
        \item 10 - 4 = 6
        \item 3 - 4 = -1
        \item 2 - 4 = -2
        \item 5 - 4 = 1
      \end{itemize}
  \end{itemize}

  \textbf{调用代码：}
  \begin{center}
  \begin{lstlisting}[linewidth=\linewidth]
result = m - 4
print(result)  # 输出：[6, -1, -2, 1]
  \end{lstlisting}
  \end{center}
  \textbf{知识点总结：}
  \begin{itemize}
    \item 可变参数\texttt{*args}的使用
    \item Python魔术方法\texttt{\_\_sub\_\_}重载运算符
    \item 列表推导式的应用
    \item 私有属性命名约定（双下划线开头）
  \end{itemize}

\end{mdframed}

\section{代码编写题}

\subsection{第1题}
编程实现，由用户输入一个不多于5位的正整数，要求求出它是几位数以及逆序打印出该数。

\begin{mdframed}
  \textbf{\color{red}【笔记】}

    \begin{lstlisting}
n = int(input())
s = str(n)  # 数字转换成字符串
print(len(s))  # 输出长度
print(s[::-1])  # 输出反转后的字符串

    \end{lstlisting}
\end{mdframed}

\subsection{第2题}
循环从用户处获得一组数据，直到用户直接输入回车退出，打印输出所有数据的和。本题不考虑输入异常情况。

\begin{mdframed}
  \textbf{\color{red}【笔记】}

    \begin{lstlisting}
n = input()
tot = 0

while n:  # 空字符串为False
    tot += eval(n)  # 转换为int类型并累加
    n = input()  # 再输入

print(tot)
    \end{lstlisting}
\end{mdframed}

\subsection{第3题}
假设os模块已导入，那么请写出C:$\backslash$Windows文件夹中（无需遍历该文件夹下的子文件夹）所有扩展名为.exe的文件

\begin{mdframed}
  \textbf{\color{red}【笔记】}

  \textbf{方法一：使用os.listdir()和字符串方法}
  \begin{lstlisting}
import os

path = r'C:\Windows'
files = os.listdir(path)

for file in files:
    if file.endswith('.exe'):
        print(file)
  \end{lstlisting}

  \textbf{方法二：使用列表推导式}
  \begin{lstlisting}
import os

path = r'C:\Windows'
exe_files = [f for f in os.listdir(path) if f.endswith('.exe')]
for file in exe_files:
    print(file)
  \end{lstlisting}

  \textbf{方法三：使用os.path.splitext()}
  \begin{lstlisting}
import os

path = r'C:\Windows'
for file in os.listdir(path):
    name, ext = os.path.splitext(file)
    if ext.lower() == '.exe':
        print(file)
  \end{lstlisting}

  \textbf{代码解释：}
  \begin{itemize}
    \item \texttt{r'C:\\Windows'}：使用原始字符串避免转义问题
    \item \texttt{os.listdir(path)}：列出指定目录下的所有文件和文件夹
    \item \texttt{file.endswith('.exe')}：判断文件名是否以.exe结尾
    \item \texttt{os.path.splitext()}：分离文件名和扩展名
    \item \texttt{ext.lower()}：转换为小写，增加匹配的容错性
  \end{itemize}

  \textbf{注意事项：}
  \begin{itemize}
    \item 此代码只列出当前目录的文件，不会递归遍历子目录
    \item 需要有访问C:\\Windows目录的权限
    \item 实际使用时可能需要异常处理
  \end{itemize}

  \textbf{os模块的常用函数介绍}

  \textbf{1. 目录操作函数：}
  \begin{itemize}
    \item \textbf{os.getcwd()}：获取当前工作目录
    \begin{lstlisting}
import os
current_dir = os.getcwd()
print(current_dir)  # 输出当前工作目录的绝对路径
    \end{lstlisting}

    \item \textbf{os.chdir(path)}：改变当前工作目录
    \begin{lstlisting}
os.chdir('/home/user/documents')  # Linux/Mac
os.chdir(r'C:\Users\Documents')   # Windows
    \end{lstlisting}

    \item \textbf{os.listdir(path)}：列出指定目录下的文件和子目录
    \begin{lstlisting}
files = os.listdir('.')  # 列出当前目录
files = os.listdir('/tmp')  # 列出/tmp目录
    \end{lstlisting}

    \item \textbf{os.mkdir(path)}：创建单级目录
    \begin{lstlisting}
os.mkdir('new_folder')  # 在当前目录创建new_folder
os.mkdir('/tmp/test')   # 在/tmp下创建test目录
    \end{lstlisting}

    \item \textbf{os.makedirs(path)}：递归创建多级目录
    \begin{lstlisting}
os.makedirs('dir1/dir2/dir3')  # 创建多级目录
os.makedirs('dir1/dir2', exist_ok=True)  # 如果存在则不报错
    \end{lstlisting}

    \item \textbf{os.rmdir(path)}：删除空目录
    \begin{lstlisting}
os.rmdir('empty_folder')  # 只能删除空目录
    \end{lstlisting}

    \item \textbf{os.removedirs(path)}：递归删除空目录
    \begin{lstlisting}
os.removedirs('dir1/dir2/dir3')  # 从内到外删除空目录
    \end{lstlisting}
  \end{itemize}

  \textbf{2. 文件操作函数：}
  \begin{itemize}
    \item \textbf{os.remove(path)}：删除文件
    \begin{lstlisting}
os.remove('test.txt')  # 删除test.txt文件
    \end{lstlisting}

    \item \textbf{os.rename(src, dst)}：重命名文件或目录
    \begin{lstlisting}
os.rename('old_name.txt', 'new_name.txt')  # 重命名文件
os.rename('old_dir', 'new_dir')  # 重命名目录
    \end{lstlisting}

    \item \textbf{os.stat(path)}：获取文件或目录的状态信息
    \begin{lstlisting}
stat_info = os.stat('file.txt')
print(stat_info.st_size)  # 文件大小（字节）
print(stat_info.st_mtime)  # 最后修改时间
    \end{lstlisting}
  \end{itemize}

  \textbf{3. os.path模块（路径操作）：}
  \begin{itemize}
    \item \textbf{os.path.exists(path)}：判断路径是否存在
    \begin{lstlisting}
if os.path.exists('file.txt'):
    print('文件存在')
    \end{lstlisting}

    \item \textbf{os.path.isfile(path)}：判断是否为文件
    \begin{lstlisting}
if os.path.isfile('test.txt'):
    print('这是一个文件')
    \end{lstlisting}

    \item \textbf{os.path.isdir(path)}：判断是否为目录
    \begin{lstlisting}
if os.path.isdir('/home'):
    print('这是一个目录')
    \end{lstlisting}

    \item \textbf{os.path.join(path1, path2, ...)}：连接路径
    \begin{lstlisting}
# 跨平台的路径连接
path = os.path.join('home', 'user', 'file.txt')
# Windows: home\user\file.txt
# Linux/Mac: home/user/file.txt
    \end{lstlisting}

    \item \textbf{os.path.split(path)}：分割路径和文件名
    \begin{lstlisting}
dir_name, file_name = os.path.split('/home/user/file.txt')
# dir_name = '/home/user'
# file_name = 'file.txt'
    \end{lstlisting}

    \item \textbf{os.path.splitext(path)}：分离文件名和扩展名
    \begin{lstlisting}
name, ext = os.path.splitext('document.pdf')
# name = 'document'
# ext = '.pdf'
    \end{lstlisting}

    \item \textbf{os.path.basename(path)}：获取文件名
    \begin{lstlisting}
file_name = os.path.basename('/home/user/file.txt')
# file_name = 'file.txt'
    \end{lstlisting}

    \item \textbf{os.path.dirname(path)}：获取目录名
    \begin{lstlisting}
dir_name = os.path.dirname('/home/user/file.txt')
# dir_name = '/home/user'
    \end{lstlisting}

    \item \textbf{os.path.abspath(path)}：获取绝对路径
    \begin{lstlisting}
abs_path = os.path.abspath('file.txt')
# 返回file.txt的绝对路径
    \end{lstlisting>

    \item \textbf{os.path.getsize(path)}：获取文件大小
    \begin{lstlisting}
size = os.path.getsize('file.txt')  # 返回字节数
    \end{lstlisting}
  \end{itemize}

  \textbf{4. 环境变量操作：}
  \begin{itemize}
    \item \textbf{os.environ}：环境变量字典
    \begin{lstlisting}
# 获取环境变量
home = os.environ.get('HOME')  # Linux/Mac
home = os.environ.get('USERPROFILE')  # Windows

# 设置环境变量
os.environ['MY_VAR'] = 'value'
    \end{lstlisting}

    \item \textbf{os.getenv(key, default=None)}：获取环境变量
    \begin{lstlisting}
path = os.getenv('PATH', '')  # 获取PATH，不存在返回空字符串
    \end{lstlisting}
  \end{itemize}

  \textbf{5. 系统相关：}
  \begin{itemize}
    \item \textbf{os.name}：操作系统名称
    \begin{lstlisting}
print(os.name)  # 'posix'(Linux/Mac) 或 'nt'(Windows)
    \end{lstlisting}

    \item \textbf{os.system(command)}：执行系统命令
    \begin{lstlisting}
os.system('ls -la')  # Linux/Mac
os.system('dir')     # Windows
    \end{lstlisting}

    \item \textbf{os.sep}：路径分隔符
    \begin{lstlisting}
print(os.sep)  # '/' (Linux/Mac) 或 '\' (Windows)
    \end{lstlisting}
  \end{itemize}

  \textbf{6. 遍历目录树 - os.walk()：}
  \begin{lstlisting}
import os

# 遍历目录树
for root, dirs, files in os.walk('/path/to/directory'):
    print(f'当前目录: {root}')
    print(f'子目录: {dirs}')
    print(f'文件: {files}')

    # 处理每个文件
    for file in files:
        file_path = os.path.join(root, file)
        print(f'文件路径: {file_path}')
  \end{lstlisting}

  \textbf{7. 实用示例：}

  \textbf{示例1：获取目录下所有特定类型的文件}
  \begin{lstlisting}
import os

def find_files(directory, extension):
    """查找目录下所有指定扩展名的文件"""
    result = []
    for root, dirs, files in os.walk(directory):
        for file in files:
            if file.endswith(extension):
                result.append(os.path.join(root, file))
    return result

# 查找所有.py文件
python_files = find_files('.', '.py')
  \end{lstlisting}

  \textbf{示例2：创建目录结构}
  \begin{lstlisting}
import os

def create_project_structure(project_name):
    """创建项目目录结构"""
    dirs = [
        f'{project_name}/src',
        f'{project_name}/tests',
        f'{project_name}/docs',
        f'{project_name}/data'
    ]

    for dir_path in dirs:
        os.makedirs(dir_path, exist_ok=True)
        print(f'创建目录: {dir_path}')
  \end{lstlisting}

  \textbf{示例3：批量重命名文件}
  \begin{lstlisting}
import os

def batch_rename(directory, old_ext, new_ext):
    """批量修改文件扩展名"""
    for filename in os.listdir(directory):
        if filename.endswith(old_ext):
            old_path = os.path.join(directory, filename)
            new_name = filename.replace(old_ext, new_ext)
            new_path = os.path.join(directory, new_name)
            os.rename(old_path, new_path)
            print(f'重命名: {filename} -> {new_name}')
  \end{lstlisting}

  \textbf{注意事项：}
  \begin{itemize}
    \item 使用os.path.join()构建路径，确保跨平台兼容性
    \item 在删除文件或目录前，先检查是否存在
    \item 处理文件操作时要考虑异常处理
    \item 使用原始字符串(r'')处理Windows路径，避免转义问题
    \item os.makedirs()的exist\_ok=True参数可以避免目录已存在的错误
  \end{itemize}
\end{mdframed}

\subsection{第4题}
在游戏应用中，经常会判断鼠标是否点击了某个人物或道具，本题将模拟此场景，编程实现判断某个点是否在某个矩形内，具体要求如下：

\textbf{(1)} 实现Point类，类中有私有属性\_\_x，\_\_y，代表鼠标的坐标，在类中实现方法get\_x，set\_x，get\_y，set\_y，使其可以分别访问相应属性。(4分)

\textbf{(2)} 实现Rectangle类，类中有私有属性x，y，代表矩形的左上角坐标，私有属性width和height，代表矩形的宽度和高度。(4分)

\textbf{(3)} 在Rectangle类中实现方法contain，判断某个点(Point类的对象)是否包含在此矩形中(4分)

\textbf{测试代码如下：}
\begin{lstlisting}
p1 = Point(20, 25)  # 点p1的坐标为x=20, y=25
# 新建矩形对象rect，其左上角坐标x=10、y=15，矩形的宽度50、高度60
rect = Rectangle(10, 15, 50, 60)
print(rect.contain(p1))

p1.set_x(80)  # 设置点p1的x坐标为80
p1.set_y(90)  # 设置点p1的y坐标为90
print(rect.contain(p1))
\end{lstlisting}

\textbf{输出结果：}
\begin{lstlisting}
True
False
\end{lstlisting}

\begin{mdframed}
  \textbf{\color{red}【笔记】}

  \textbf{完整代码实现：}
  \begin{lstlisting}
class Point:
    def __init__(self, x, y):
        self.__x = x
        self.__y = y

    def set_x(self, x):
        self.__x = x

    def set_y(self, y):
        self.__y = y

    def get_x(self):
        return self.__x

    def get_y(self):
        return self.__y

class Rectangle:
    def __init__(self, x, y, width, height):
        self.x = x
        self.y = y
        self.width = width
        self.height = height

    def contain(self, point):
        if self.x <= point.get_x() <= self.x + self.width \
        and self.y <= point.get_y() <= self.y + self.height:
            return True
        return False
  \end{lstlisting}

  \textbf{Python类详解：}

  \textbf{1. 类的基本概念：}
  \begin{itemize}
    \item 类（Class）是对象的模板，定义了对象的属性和方法
    \item 对象（Object）是类的实例，具有类定义的属性和行为
    \item 使用关键字\texttt{class}定义类，类名通常采用驼峰命名法
  \end{itemize}

  \textbf{2. 构造函数\_\_init\_\_：}
  \begin{itemize}
    \item \texttt{\_\_init\_\_}是特殊方法，用于初始化对象
    \item 创建对象时自动调用，相当于其他语言的构造函数
    \item \texttt{self}参数代表对象本身，必须是第一个参数
    \item 可以接收参数来初始化对象的属性
  \end{itemize}

  \textbf{3. 属性访问控制：}
  \begin{lstlisting}
# 公有属性：可以直接访问
self.x = x          # 公有属性

# 私有属性：不能直接从外部访问
self.__x = x        # 私有属性（双下划线开头）

# 访问示例
p = Point(10, 20)
print(p.x)          # 如果x是公有属性，可以直接访问
# print(p.__x)      # 错误！私有属性不能直接访问
print(p.get_x())    # 通过getter方法访问私有属性
  \end{lstlisting}

  \textbf{4. Getter和Setter方法：}
  \begin{itemize}
    \item Getter方法：用于获取私有属性的值
    \item Setter方法：用于设置私有属性的值
    \item 提供了对私有属性的受控访问
    \item 可以在方法中添加验证逻辑
  \end{itemize}

  \textbf{5. 方法定义和调用：}
  \begin{lstlisting}
# 方法定义
def method_name(self, parameters):
    # 方法体
    return value

# 方法调用
obj = ClassName()
result = obj.method_name(arguments)
  \end{lstlisting}

  \textbf{6. 代码逻辑分析：}

  \textbf{Point类：}
  \begin{itemize}
    \item 私有属性\texttt{\_\_x}和\texttt{\_\_y}存储坐标
    \item 提供getter和setter方法访问坐标值
    \item 封装了坐标点的基本操作
  \end{itemize}

  \textbf{Rectangle类：}
  \begin{itemize}
    \item 公有属性存储矩形的位置和尺寸信息
    \item \texttt{contain}方法判断点是否在矩形内部
    \item 判断条件：\texttt{x <= point\_x <= x + width}且\texttt{y <= point\_y <= y + height}
  \end{itemize}

  \textbf{7. 面向对象编程优势：}
  \begin{itemize}
    \item \textbf{封装}：将数据和操作封装在类中
    \item \textbf{抽象}：隐藏实现细节，只暴露必要接口
    \item \textbf{重用性}：类可以被多次实例化使用
    \item \textbf{维护性}：便于代码的维护和扩展
  \end{itemize}

  \textbf{8. 类装饰器详解：}

  \textbf{8.1 @property装饰器：}
  \begin{itemize}
    \item @property将方法转换为属性，提供更优雅的访问方式
    \item 可以像访问属性一样调用方法，无需加括号
    \item 提供了getter、setter、deleter的完整支持
    \item 比传统的get/set方法更加Pythonic
  \end{itemize}

  \textbf{8.2 使用@property重写Point类：}
  \begin{lstlisting}
class Point:
    def __init__(self, x, y):
        self.__x = x
        self.__y = y

    @property
    def x(self):
        """获取x坐标"""
        return self.__x

    @x.setter
    def x(self, value):
        """设置x坐标，可以添加验证"""
        if not isinstance(value, (int, float)):
            raise TypeError("坐标必须是数字")
        self.__x = value

    @x.deleter
    def x(self):
        """删除x坐标"""
        print("删除x坐标")
        del self.__x

    @property
    def y(self):
        """获取y坐标"""
        return self.__y

    @y.setter
    def y(self, value):
        """设置y坐标"""
        if not isinstance(value, (int, float)):
            raise TypeError("坐标必须是数字")
        self.__y = value

    @property
    def distance_from_origin(self):
        """计算到原点的距离（只读属性）"""
        return (self.__x ** 2 + self.__y ** 2) ** 0.5
  \end{lstlisting}

  \textbf{8.3 @property的使用示例：}
  \begin{lstlisting}
# 创建点对象
p = Point(3, 4)

# 使用@property装饰的方法，像属性一样访问
print(p.x)                    # 输出: 3 (调用getter)
print(p.y)                    # 输出: 4 (调用getter)
print(p.distance_from_origin) # 输出: 5.0 (只读属性)

# 使用setter设置值
p.x = 10                      # 调用setter
p.y = 20                      # 调用setter

# 验证功能
try:
    p.x = "invalid"           # 触发TypeError异常
except TypeError as e:
    print(e)                  # 输出: 坐标必须是数字

# 删除属性
del p.x                       # 调用deleter
  \end{lstlisting}

  \textbf{8.4 传统方法vs @property对比：}
  \begin{lstlisting}
# 传统getter/setter方法
point1 = Point(10, 20)
print(point1.get_x())         # 需要调用方法
point1.set_x(30)              # 需要调用方法

# 使用@property装饰器
point2 = Point(10, 20)
print(point2.x)               # 像访问属性一样
point2.x = 30                 # 像设置属性一样
  \end{lstlisting}

  \textbf{8.5 @property的优势：}
  \begin{itemize}
    \item \textbf{简洁性}：访问语法更加直观，无需get/set前缀
    \item \textbf{封装性}：可以在getter/setter中添加验证逻辑
    \item \textbf{兼容性}：可以将普通属性升级为property而不破坏接口
    \item \textbf{只读属性}：可以创建计算属性（如distance\_from\_origin）
    \item \textbf{延迟计算}：属性值可以在访问时动态计算
  \end{itemize}

  \textbf{8.6 其他常用装饰器：}
  \begin{lstlisting}
class MyClass:
    @staticmethod
    def static_method():
        """静态方法，不需要self参数，不依赖实例"""
        return "这是静态方法"

    @classmethod
    def class_method(cls):
        """类方法，接收cls参数，可以访问类属性"""
        return f"这是{cls.__name__}的类方法"

    @property
    def read_only_prop(self):
        """只读属性"""
        return "只读值"

# 使用示例
obj = MyClass()
print(MyClass.static_method())    # 直接通过类调用
print(MyClass.class_method())     # 通过类调用类方法
print(obj.read_only_prop)         # 访问只读属性
  \end{lstlisting}

  \textbf{8.7 实际应用场景：}
  \begin{itemize}
    \item \textbf{数据验证}：在setter中验证输入数据的有效性
    \item \textbf{计算属性}：根据其他属性动态计算值
    \item \textbf{缓存机制}：在属性中实现计算结果的缓存
    \item \textbf{兼容性}：将直接属性访问升级为方法调用
    \item \textbf{日志记录}：在属性访问时记录日志
  \end{itemize}

  \textbf{9. 扩展知识：}
  \begin{lstlisting}
# 类属性和实例属性
class MyClass:
    class_var = "类属性"    # 类属性，所有实例共享

    def __init__(self, value):
        self.instance_var = value  # 实例属性，每个实例独有

# 特殊方法（魔术方法）
class Point:
    def __str__(self):
        return f"Point({self.__x}, {self.__y})"

    def __repr__(self):
        return f"Point({self.__x}, {self.__y})"

    def __eq__(self, other):
        return self.__x == other.__x and self.__y == other.__y
  \end{lstlisting}

\end{mdframed}


\end{document}